\documentclass[10pt,a4paper,ssfamily]{exam}
\usepackage[brazil]{babel}
\usepackage[numbers,sort&compress]{natbib}
\usepackage[utf8]{inputenc}
\renewcommand{\baselinestretch}{1.0}
\usepackage{epsfig}
\usepackage{mhchem}
\usepackage{graphicx}
\usepackage{makeidx}
\usepackage{multicol}
\usepackage{multirow}
\usepackage{amssymb}
\usepackage{amsmath}
\DeclareMathOperator{\sen}{sen}
\newcommand{\listfont}{\bf\fontencoding{OT1}\sffamily\large}
\newcommand{\titlesfont}{\fontencoding{OT1}\sffamily\large}
\renewcommand{\baselinestretch}{1.0}
\newcommand{\Mg}{Mg$^{2+}$}
\renewcommand{\t}{\textrm}
\newcommand{\cc}[1]{[\textrm{#1}]}
\usepackage{enumitem}
\setlist{nosep, topsep=2pt, partopsep=0pt, parsep=0pt, itemsep=2pt}
\newcommand{\bmat}{\left[\begin{array}{cc}}
\newcommand{\emat}{\end{array}\right]}
\newcommand{\nomera}{
  \noindent
  \begin{tabular}{|p{0.7\textwidth}|p{0.3\textwidth}|}
  Nome: & RA: \\
  \hline
  \end{tabular}
}

\newcommand{\qini}{\begin{enumerate}[resume]\item}
\newcommand{\qend}{\end{enumerate}}

\newcommand{\oC}{$^\circ$C}
\newcommand{\kcalmol}{~kcal~mol$^{-1}$}
\newcommand{\moll}{~mol~L$^{-1}$}
\newcommand{\molkg}{~mol~kg$^{-1}$}
\newcommand{\gl}{~g~L$^{-1}$}
\newcommand{\e}[1]{$\times 10^{#1}$}

\newcommand{\eps}{\varepsilon}

\newcommand{\Resposta}[1]{\vspace{0.5cm}\noindent{\bf Resposta:}\\ #1
\begin{center}\rule{\textwidth}{0.4pt}\end{center}}
\renewcommand{\Resposta}[1]{}
\newcommand{\resposta}[1]{\vspace{0.5cm}\noindent{\bf Resposta:}\\ #1
\begin{center}\rule{\textwidth}{0.4pt}\end{center}}


\begin{document}
\sffamily
\pagestyle{empty}
\nomera

\begin{center}
{\bf QG101 - Química Geral}\\
{\bf Aula 6 - Ligações Iônicas}\\
\underline{Justifique suas respostas.}
\end{center}

\vspace{0.2cm}
\begin{center}{\bf Questões}\end{center}

% ---------------------------------------------------------------
\qini
O \textbf{cloreto de potássio} (\ce{KCl}) e o \textbf{óxido de
magnésio} (\ce{MgO}) são dois sólidos iônicos com a mesma estrutura
cristalina (tipo \ce{NaCl}).

\begin{enumerate}
  \item Escreva as configurações eletrônicas dos íons
        \ce{K+}, \ce{Cl-}, \ce{Mg^{2+}} e \ce{O^{2-}}.
        Identifique com quais gases nobres eles são isoeletrônicos.

  \item A equação de Born--Landé para a energia de rede é:
        \[
          U = -\frac{N_A\, A\, |z_+||z_-|\, e^2}{4\pi\varepsilon_0\, r_0}
              \!\left(1 - \frac{1}{n}\right)
        \]
        Usando \emph{apenas} os parâmetros $|z_+||z_-|$ e $r_0$
        ($r_0(\text{KCl}) = 319$~pm; $r_0(\text{MgO}) = 211$~pm),
        explique \textbf{qualitativamente} por que
        $|U(\ce{MgO})| \gg |U(\ce{KCl})|$.

  \item O ponto de fusão do \ce{MgO} é $2852$\oC{} e o do \ce{KCl}
        é $770$\oC. Relacione essa diferença à energia de rede.
\end{enumerate}

\vspace{0.3cm}
\qend

% ---------------------------------------------------------------
\qini
\textbf{Constante de Madelung unidimensional.}
Considere uma cadeia linear infinita de íons com cargas alternadas
$+e$ e $-e$, separados por distância $d$:

\[
  \cdots\; +\; -\; +\; -\; +\; -\; +\;\cdots
\]

\begin{enumerate}
  \item Tomando um íon de referência $+$ na origem,
        mostre que a energia eletrostática por par é proporcional
        à série:
        \[
          S = 1 - \tfrac{1}{2} + \tfrac{1}{3} - \tfrac{1}{4}
              + \cdots = \sum_{n=1}^{\infty}\frac{(-1)^{n+1}}{n}
        \]

  \item Sabendo que $S = \ln 2$, calcule a constante de Madelung
        unidimensional $A_{1D} = 2\ln 2$ e compare com
        $A_{3D} = 1{,}748$ (estrutura \ce{NaCl}).
        Por que a constante tridimensional é maior?
\end{enumerate}

\vspace{0.3cm}
\qend

% ---------------------------------------------------------------
\qini
\textbf{Ciclo de Born--Haber para o \ce{KCl}.}
Use os dados abaixo para calcular a energia de rede $U$ do \ce{KCl}.

\begin{center}
\begin{tabular}{lc}
\hline
Grandeza & $\Delta H$ (kJ\,mol$^{-1}$)\\
\hline
$\Delta H_f^\circ$(\ce{KCl}, s) & $-437$ \\
$\Delta H_{\text{sub}}$(\ce{K}) & $+89$ \\
$\frac{1}{2}D$(\ce{Cl2}) & $+121$ \\
$\text{EI}_1$(\ce{K}) & $+419$ \\
$\text{AE}$(\ce{Cl}) & $+349$ \\
\hline
\end{tabular}
\end{center}

\begin{enumerate}
  \item Desenhe o ciclo de Born--Haber indicando cada etapa e
        sua variação de entalpia.

  \item Aplique a Lei de Hess para obter a expressão de $U$ em
        função das grandezas tabeladas e calcule o valor numérico.

  \item O valor experimental de $U(\ce{KCl})$ é $-717$~kJ/mol.
        Compare com o seu resultado e comente.
\end{enumerate}

\vspace{0.3cm}
\qend

\end{document}
